\section{Introduction}\label{sec:introduction}

A sentence with one existential and one universal real block asks whether
there is a real tuple for which every real tuple in a second block satisfies
specified polynomial conditions. We study the discrete language
$\EATR$ of true sentences
\begin{equation}\label{eq:input-sentence}
 \exists x\in\R^r\ \forall y\in\R^s:\ \Phi(x,y).
\end{equation}
Here $r,s\ge0$ are represented by explicit lists of variables, and $\Phi$
is an explicitly encoded Boolean formula using AND, OR, and NOT over atoms
$p=0$ and $p>0$. Each $p$ is given by an explicitly listed acyclic arithmetic
circuit with signed binary integer constants and gates $+,-,\times$.
Predecessor indices specify the gates; sharing is allowed, and there are
no division or exponentiation gates. All operations and comparisons in
\eqref{eq:input-sentence} are interpreted exactly. The input length $N$
is the total binary encoding length, and malformed encodings are rejected.

All computational bounds below are uniform bit-time bounds. Let $\PP$
denote polynomial-time probabilistic computation with fair random bits,
accepting exactly the yes-instances with probability greater than $1/2$.
The counting hierarchy, introduced by Wagner~\cite{Wagner1986}, is defined by
\[
 \Clevel0=\Pclass,\qquad
 \Clevel{j+1}=\bigcup_{A\in\Clevel j}\PP^A,
 \qquad \CH=\bigcup_{j\ge0}\Clevel j.
\]
Membership uses a fixed oracle language and a fixed uniform machine.

\begin{theorem}\label{thm:main}
There is an absolute integer $j$ such that
\[
                         \EATR\in\Clevel j.
\]
The level $j$ is independent of the input length, the numbers of variables,
the degrees of the circuit polynomials, and their coefficient magnitudes.
\end{theorem}

The existential theory of the reals, denoted $\ETR$, is the same decision
problem with only the existential block. Write $\existsR$ for the class
of languages polynomial-time many-one reducible to $\ETR$. Finer
accounting for this fragment gives $\ETR\in\Clevel{26}$ and hence
$\existsR\subseteq\Clevel{26}$; Theorem~\ref{thm:etr-bound} proves this
explicit bound. The main result allows a universal block of unrestricted
dimension as well.

\subsection{History and significance}

Tarski proved decidability of the first-order theory of real closed
fields~\cite{Tarski1951}. Collins's cylindrical algebraic decomposition
gave a geometric method for eliminating quantifiers~\cite{Collins1975},
and Canny established the polynomial-space upper bound for the existential
fragment~\cite{Canny1988}. For each fixed number of real quantifier blocks,
Renegar's decision bounds yield a polynomial-space algorithm
\cite[Theorem~1.1]{Renegar1992II}. This applies to the circuit encoding
above: the quartic reduction in Lemma~\ref{lem:quartic-fibers} preserves
the two quantifier blocks. Our theorem places this two-block language
in a fixed level of the counting hierarchy.
Real quantifier hierarchies distinguish the
number and order of blocks; Schaefer and \v{S}tefankovi\v{c} develop their
finite-encoding formulation~\cite{SchaeferStefankovic2025}.
The existential class already captures many geometric realization and
continuous feasibility problems; Schaefer, Cardinal, and Miltzow survey
these connections~\cite{SchaeferCardinalMiltzow2026}.

The counting hierarchy gives a way to decide some exact algebraic
questions without constructing their enormous intermediate integers.
Allender, B\"urgisser, Kjeldgaard-Pedersen, and Miltersen proved that
$\mathsf{PosSLP}$---positivity of the integer output of a division-free
straight-line program---belongs to $\CH$
\cite{AllenderBurgisserKjeldgaardPedersenMiltersen2009}.
Their modular-product and approximate Chinese-remainder techniques are
the counting tools behind our treatment of implicit polynomial families.
Every fixed counting level lies in $\PSPACE$: a depth-first simulation
enumerates random choices and answers the fixed-depth oracle queries
using polynomial space.

Recent counting-hierarchy bounds address other polynomial feasibility
problems. Andrews, Garg, and Schost place Hilbert's Nullstellensatz over
$\Q$ and several other fields in $\CH$~\cite{AndrewsGargSchost2026};
the common zeros there lie in an algebraic closure.
Balaji, Shirmohammadi, Tavenas, and Worrell prove counting-hierarchy
bounds for approximate polynomial satisfiability over $\Q$ and finite
fields~\cite{BalajiShirmohammadiTavenasWorrell2026}, where the question
is membership of zero in the Zariski closure of a polynomial image
over the algebraic closure.
Exact real quantification must also control the reality of solutions,
their attainment, and the dependence of a quantified fiber on its
parameters. These are the issues addressed by the argument here.

Our algebraic tools belong to the critical-point and deformation methods
of real algebraic geometry~\cite{BasuPollackRoy2006}. Multiplication
determinants are evaluated through the classical relation between
residues and traces~\cite{SchejaStorch1975,CattaniDickensteinSturmfels1996}.
We prove the needed identities and computational rules directly, including
their behavior under parameter specialization. These computational rules
apply to implicitly represented families that need not have short
coefficient lists or arithmetic circuits.

\subsection{From quantified fibers to short algebraic witnesses}

Fix the existential tuple $x$. Introducing gate and auxiliary variables
turns failure of $\Phi(x,y)$ into a real zero of a nonnegative quartic
$F(x,z)$, where $z$ includes $y$ and the auxiliaries. Thus we seek an
$x$ for which the fiber $F(x,\cdot)$ has no real zero. Call such an $x$
\emph{good}. Merely testing fibers does not provide a finite description
of a good parameter. The main geometric argument produces such descriptions.

For $u\ge0$, consider the minimum
\[
 m_x(u)=\min_{z\in\R^n}\left(\sum_{i=1}^n z_i^6+6uF(x,z)\right).
\]
The sixth powers keep bounded sublevel sets compact. It follows that
$m_x(u)$ remains bounded as $u\to\infty$ precisely when the fiber has
a zero. At a minimizing point the critical equations have pure fifth
powers as their leading terms. Their quotient algebra has dimension
$D=5^n$, and multiplication by the minimized polynomial has a monic
characteristic polynomial $Q(x,u,Y)$ satisfying
$Q(x,u,m_x(u))=0$.

This identity supplies a uniform growth test: $x$ is good precisely when
$m_x(v^{D+1})>v$ for all sufficiently large $v$. Form
\[
                         E(x,v)=Q(x,v^{D+1},v).
\]
The coefficient of $v^D$ is always $1$, so $E(x,\cdot)$ is never zero.
Partition parameter space into the sets on which $E(x,\cdot)$ has each
possible degree. On each such degree piece, polynomial root bounds give
a locally uniform threshold beyond which $E$ does not vanish.
Continuity of the minimum then makes goodness locally constant on that
piece. Introducing the reciprocal of its leading coefficient realizes
the piece as a closed algebraic set. Both the good and nongood parts
of this lift are closed.

A second minimization argument now extracts candidate coordinates.
Choose a compact sublevel set of a coercive polynomial with a good
point in its interior. Within this set, the good and nongood parts of
the lift are separated. Minimize the coercive polynomial plus a growing
nonnegative polynomial penalty that vanishes exactly on the lift,
restricting to a compact neighborhood that avoids the nongood part. A subsequence of
minimizers converges to a good point and is eventually interior to the
neighborhood, so these minimizers satisfy the unconstrained critical
equations. The highest parameter coefficients of their coordinate
characteristic polynomials therefore vanish at the limit coordinates.
This produces finitely indexed univariate polynomial families whose roots
contain the coordinates of a suitable witness, without computing the good part or its connected
components. Section~\ref{sec:two-blocks} proves this geometric
construction and completes the main theorem.

Two computational ingredients make these candidate coordinates usable.
First, all polynomials just described have degree and logarithmic
coefficient bounds with polynomial-length binary encodings, and their
finite-field evaluations lie at fixed counting levels. We call such
families \emph{controlled}. Sections~\ref{sec:counting} and
\ref{sec:algebra} establish their closure rules and the evaluation of
multiplication determinants without constructing the matrices.

Second, a chosen real root needs a short description even when its
polynomial has exponential degree. Derivative signs distinguish real
roots, as in Thom encodings~\cite{BasuPollackRoy2006}, but the full sign
string can be too long. Section~\ref{sec:labels} compresses that string
to a weighted integer sum with weights computed modulo a short prime,
and realizes the resulting
label by a polynomial positivity test. Every desired real root can
be isolated, including a root of a polynomial with multiplicities.
The sign approximation uses geometric products of the type studied
by Beigel, Reingold, and Spielman~\cite{BeigelReingoldSpielman1995}.
This arbitrary-root labeling rule and the coordinate construction
are reusable ingredients of the proof.

Finally, Section~\ref{sec:existential} proves an existential decision
procedure for controlled families. For a guessed coordinate description,
we use it twice: the described parameter set must be nonempty, and
it must contain no parameter with a zero of $F$. A suitable choice
of root labels singles out the good point constructed above. All
choices have polynomial total bit length, and all family constructions
have fixed depth. Counting over these choices therefore gives one
fixed level of the hierarchy.
